\section*{Chapter \ref{ch:g2:born-grow-die} --- Born, Growing, Dying}

\begin{solution}{exo:g2:born-grow-die:1}
Born, young, \omterm{def:g2:born-grow-die:cycle}{adult}, \omterm{def:g2:born-grow-die:cycle}{old age} (ending in death).
\end{solution}

\begin{solution}{exo:g2:born-grow-die:2}
At the \omterm{def:g2:born-grow-die:cycle}{adult} \omterm{def:g2:born-grow-die:cycle}{stage}.
\end{solution}

\begin{solution}{exo:g2:born-grow-die:3}
Oldest to youngest: the grandmother, the mother, the schoolchild,
the baby.
\end{solution}

\begin{solution}{exo:g2:born-grow-die:4}
A cherry seed sprouts (\omterm{def:g2:born-grow-die:cycle}{birth}); the young tree grows taller for
years; the \omterm{def:g2:born-grow-die:cycle}{adult} tree makes blossom and cherries --- and so new
seeds; then it grows old and dies.
\end{solution}

\begin{solution}{exo:g2:born-grow-die:5}
For example: a baby gets teeth, loses them and gets stronger ones;
a chick's yellow fluff becomes real \omterm{def:g1:animals-around-us:covering}{feathers}; a young tree's soft
green \omterm{def:g1:plants-around-us:parts}{stem} turns to hard wood.
\end{solution}

\begin{solution}{exo:g2:born-grow-die:6}
Shortest to longest: the butterfly (a few weeks), the cat (ten to
twenty years), the oak tree (hundreds of years).
\end{solution}

\begin{solution}{exo:g2:born-grow-die:7}
Because before the road ends, \omterm{def:g2:born-grow-die:cycle}{adults} have young, and the young
begin the road again --- that is the purple arrow looping back.
One life is a road; life over the generations turns like a wheel.
\end{solution}

\begin{solution}{exo:g2:born-grow-die:8}
Open answer. The photographs usually show baby, child, \omterm{def:g2:born-grow-die:cycle}{adult} ---
\omterm{def:g2:born-grow-die:cycle}{old age} may be the \omterm{def:g2:born-grow-die:cycle}{stage} still to come.
\end{solution}

\begin{solution}{exo:g2:born-grow-die:9}
The old tree's apples carry pips --- seeds --- and each planted
pip is a young apple tree: the next generation. When the old tree
dies, its young continue the cycle, so the garden keeps its apple
trees.
\end{solution}

\begin{solution}{exo:g2:born-grow-die:10}
No --- it is because the pebble was \omterm{def:g1:living-or-not:nonliving}{never alive}. Only \omterm{def:g1:living-or-not:living}{living
things} are born, grow and die; a pebble has no \omterm{def:g2:born-grow-die:cycle}{life cycle} at all,
so ``never dying'' is not strength, it is not living.
\end{solution}
